\begin{textsample}{POS Dim 2 – rn – Score 20.00 – rn/rn\_inf\_40.txt}  \label{ex:f2_pos_209}
9. \textbf{GESTÃO} E COORDENAÇÃO DO TRABALHO EDUCATIVO A Educação Infantil, primeira etapa da Educação Básica, é oferecida em creches e pré-escolas, as quais se caracterizam como espaços institucionais não domésticos que constituem estabelecimentos educacionais públicos ou privados que educam e cuidam de crianças de 0 a 5 anos de idade no período diurno, em jornada integral ou parcial, regulados e supervisionados por órgão competente do sistema de ensino e submetidos a controle social. ( BRASIL, 2009a, p. 01 ).
A Educação Infantil como etapa educativa e espaço institucional que deve possibilitar a aprendizagem e o desenvolvimento de crianças, envolve vários sujeitos e âmbitos da \textbf{gestão}, coordenação, supervisão e articulação na organização de suas práticas pedagógicas.
Neste \textbf{capítulo} destacam -se algumas orientações breves, relativas a essa organização das redes de ensino e das \textbf{instituições} de Educação Infantil, construídas em diálogo com professores e demais profissionais que contribuíram na consulta pública.
De início, é importante destacar que o \textbf{currículo} da/para educação infantil também envolve a \textbf{gestão} e a coordenação do trabalho educativo, bem como, se orienta e se desenvolve considerando os modos como as redes públicas e privadas se organizam para esse atendimento educacional.
Por isso, considerou -se importante abordar este tema neste documento curricular.
APROFUNDANDO O TEMA \textbf{Regulamentação} das \textbf{Instituições} públicas e privadas de Educação Infantil Os sistemas de ensino têm autonomia para complementar a \textbf{legislação} nacional por meio de normas próprias, específicas e adequadas às características locais.
O Conselho Estadual de Educação do Rio Grande do Norte - CEE, na Resolução 01/2016, fixa as normas para a organização e o funcionamento da Educação Infantil no Estado.
Os \textbf{municípios} que possuem sistema de ensino devem dispor de \textbf{regulamentação} e normas que orientam a criação, a autorização, o funcionamento, a supervisão e a avaliação das \textbf{instituições} de Educação Infantil, sendo essa \textbf{regulamentação} de competência do Conselho Municipal de Educação - CME.
Os demais \textbf{municípios} \textbf{permanecem} integrados ao sistema estadual e seguem as normas definidas pelo CEE.
De modo geral, as normas abordam critérios e exigências que balizam o funcionamento das \textbf{instituições} de educação infantil, tais como : formação dos professores ; espaços físicos, incluindo parâmetros para assegurar higiene, segurança, conforto ; número de crianças por professor ; proposta pedagógica ; \textbf{gestão} dos estabelecimentos e \textbf{documentação} exigida.
O atendimento na educação infantil deve, portanto, observar leis e normas \textbf{municipais}, estaduais e \textbf{federais} ( aportes legais apresentados no primeiro \textbf{capítulo} deste documento ) referentes a esta etapa educativa.
Assim como, a Lei Orgânica Municipal, as exigências referentes à Construção Civil e ao Código Sanitário, entre outros.
A adequada organização e estruturação do sistema de ensino é essencial para que a educação infantil se efetive como \textbf{política} educacional.
Não basta o Conselho definir as normas, é preciso que a Secretaria de Educação oriente as \textbf{instituições} e dê os suportes \textbf{técnico}, pedagógico e financeiro necessários para que elas consigam se adequar às exigências da \textbf{regulamentação}.
As \textbf{instituições} de Educação Infantil, por sua vez, devem promover as devidas adequações às regras do respectivo sistema de ensino.
Uma estratégia importante para garantir um atendimento de qualidade é o desenvolvimento de Políticas \textbf{municipais} de Educação infantil, que, por sua vez, além dos aspectos aqui destacados para \textbf{normatização}, devem considerar as metas do PNE ( BRASIL, 2014 ), como, por exemplo, a \textbf{garantia} para que as crianças de quatro e cinco anos estejam na Pré-Escola e seja ampliada a \textbf{oferta} em Creches, de forma que \textbf{atendam}, no mínimo, 50 \% das crianças de até três anos, entre outras. ( BRASIL, 2013 )

% matched lemmas: atender, capítulo, currículo, documentação, federal, garantia, gestão, instituição, legislação, municipal, município, normatização, oferta, permanecer, política, regulamentação, técnico
\end{textsample}
